% "{'chapitre':'slci_multiphy','classe':('PSI'),'type':('application'),'titre':'Mise à l\\'eau d\\'un robot sous-marin','source':'Concours Centrale - MP 2019','comp':['B2-02',],'corrige':True}"

%\setchapterimage{fig_00.jpg}
\chapter*{Application \arabic{cptApplication} \\ 
Véhicule hybride TOYOTA PRIUS -- \ifprof Corrigé \else Sujet \fi}
\addcontentsline{toc}{section}{Application \arabic{cptApplication} : Véhicule hybride TOYOTA PRIUS -- \ifprof Corrigé \else Sujet \fi}

\iflivret \stepcounter{cptApplication} \else
\ifprof  \stepcounter{cptApplication} \else \fi
\fi

\setcounter{question}{0}
\marginnote{Concours CCS -- PSI}





\subsection*{Question de cours}
\question{Donner la strucutre de la chaîne de puissance. On précisera en particulier le nom de chacune des fonctions.}

\question{Donner 3 exemples de composants permettant de réaliser chacune des fonctions.}

\subsection*{Véhicule hybride TOYOTA PRIUS}
\marginnote{D'après ressources de Florestan Mathurin.}

\begin{marginfigure}
\includegraphics[width=\linewidth]{fig_01}
\end{marginfigure}
Dans le contexte actuel d’économie des énergies fossiles et de
réduction des émissions de gaz nocifs, le système de propulsion
hybride constitue une alternative intéressante à la propulsion
classique par moteur thermique seul car il permet de réduire la
consommation.

La spécificité de la solution retenue sur la Prius consiste à :
\begin{itemize}
\item récupérer l’énergie du véhicule lors du freinage;
\item exploiter le moteur thermique à son rendement optimal.
\end{itemize}

La technologie hybride de TOYOTA, nommée HSD (Hybrid
Synergy Drive) associe un moteur thermique à essence et
sa transmission, à deux machines électriques et une
batterie de puissance.

Le schéma de principe ci-contre met en évidence les deux
machines électriques (le moteur électrique et la
génératrice) reliées au moteur thermique par un réducteur
à train épicycloïdal.

A partir de la position de la pédale d’accélérateur et de la
vitesse du véhicule, le calculateur détermine la vitesse de
rotation optimale du moteur thermique et la consigne
d’ouverture du papillon des gaz.

La puissance en sortie du moteur thermique est transmise, grâce à un train épicycloïdal, à la chaîne silencieuse
et à la génératrice. Un asservissement en vitesse de la génératrice permet de contrôler la vitesse de rotation
du moteur thermique.

Le répartiteur de puissance gère les échanges de puissance électrique entre la génératrice, le moteur
électrique et la batterie. Le moteur électrique entraîne la chaîne silencieuse, seul ou en complément du
moteur thermique. Il récupère également l’énergie cinétique ou potentielle du véhicule lors des phases de
ralentissement.

\begin{center}
\includegraphics[width=\linewidth]{fig_02}
\end{center}

\question{Compléter le diagramme chaîne d’information / chaîne de puissance.}

\question{Dans la chaîne de puissance, repasser en rouge sur le diagramme les flèches représentant les flux
correspondant à une puissance mécanique et en rose les flèches représentant les flux correspondant à une
puissance électrique.}

\begin{center}
\includegraphics[width=\linewidth]{fig_03}
\end{center}


